% Chapter 8 of Optimizers in the Wild: the edge of stability.
\newcommand{\eosHmin}{\cH_{\min}}
\newcommand{\eosArea}{\mathcal{J}}

\casechapter{The Edge of Stability}{Thresholds, self-stabilization and central flows}{Learning rates at the stability limit}{ch:eos}
\begin{refsection}[chapters/08-eos.bib]
\nocite{HydrogenAtoms2026}

\begin{chaptersummary}
Learning rates are chosen as large as training tolerates, and full-batch training then runs with its curvature pinned at the stability threshold of the optimizer: the sharpness at $2/\eta$ for gradient descent, the preconditioned sharpness at $38/\eta$ for Adam with $\beta_1=0.9$. We derive these thresholds from one stability test, and then study the cubic model of Damian, Nichani and Lee, in which a single direction oscillates and the oscillation pushes the curvature down. In its continuous limit we prove that every cycle of oscillation conserves an energy and closes, and that over each cycle the curvature averages exactly $2/\eta$ while the squared oscillation averages exactly twice the rate of progressive sharpening; for the gradient-descent map itself we prove that the same energy increases at every step, so that the cycles spiral slowly outward, and we derive the rate. Averaging over the cycles gives the central flow, gradient flow projected onto the region where the step size is stable. The model takes progressive sharpening as an input and does not explain it.
\end{chaptersummary}

\section{In the wild: the largest learning rate that does not diverge}\label{sec:eos:wild}

Gradient descent on a quadratic multiplies the component of $\theta_t-\theta^\star$ along an eigenvector of the Hessian with eigenvalue $\lambda_i$ by $1-\eta\lambda_i$ at every step, so it converges from every initialization exactly when $\eta\lambda_i<2$ for every eigenvalue (\cref{ch:prologue}). The largest eigenvalue of the Hessian, $S(\theta)=\lmax(\nabla^2\Loss(\theta))$, is the \term{sharpness}\index{sharpness}, and $\eta S<2$ is the classical stability condition. In this chapter the learning rate and the step size are the same number $\eta$; we say learning rate when we describe practice and step size in the theory. Read as advice, the stability condition says: estimate the curvature, then choose the step. Networks are not trained that way.

\begin{practice}[folklore]
Choose the learning rate by a sweep on a logarithmic grid, with neighbouring values a factor of two or three apart, and keep the value with the lowest loss after a fixed budget; it is usually among the largest values that train without diverging, and a run that diverges is restarted with a smaller value. The curvature of the loss is not measured, neither at initialization nor later. Momentum, and Adam's average of the gradient with $\beta_1=0.9$ \parencite{KingmaBa2015}, stay at their defaults while the learning rate is swept.
\end{practice}

The rule is folklore in the sense of this book: it describes how sweeps are commonly run, not a measured regularity, and the \emph{Deep Learning Tuning Playbook} \parencite{GodboleEtAl2023} is the most systematic guide to running them. What the rule leaves unsaid is that the curvature is not a fixed property of the loss that the learning rate has to respect. It responds to the learning rate.\sidenote{That the step size limits where gradient descent can be stable, and so selects flat minima, is the dynamical-stability argument of \textcite{WuMaE2018}; \textcite{JastrzebskiEtAl2019} observed that the sharpness reached early in training adapts to the step size.}

\begin{observation}[Progressive sharpening and the edge of stability; \textcite{CohenEtAl2021}]\label[observation]{obs:eos:gd}
Train a network by full-batch gradient descent with a fixed step size $\eta$. (i) While $S(\theta_t)<2/\eta$, the sharpness tends to rise. (ii) Once it reaches $2/\eta$, it stops rising and hovers at or just above $2/\eta$, while the training loss keeps decreasing over long time scales, though not monotonically over short ones. (iii) With heavy-ball momentum $\beta$ the sharpness settles instead near $(2+2\beta)/\eta$, and with Nesterov momentum near $(2+2\beta)/((1+2\beta)\eta)$: in each case the largest sharpness at which the method is stable on a quadratic.
\end{observation}

Cohen and coauthors named the first stage \term{progressive sharpening}\index{progressive sharpening} and the second the \term{edge of stability}\index{edge of stability}, and documented both for fully connected, convolutional and transformer networks, with the square loss and with the cross-entropy loss. For momentum they derived the thresholds in (iii), which they call maximal stable sharpnesses, and observed the sharpness rise to them and level off in runs with $\eta=0.01$. The observation is robust for full-batch gradient descent, and its scope is limited in two ways that will matter later: progressive sharpening is weaker for wide or shallow networks and for easy tasks, and with the cross-entropy loss the sharpness typically falls at the end of training \parencite{CohenEtAl2021}.

\begin{observation}[The adaptive edge of stability; \textcite{CohenEtAl2022}]\label[observation]{obs:eos:adam}
Train a network full batch with an adaptive method, such as Adam or RMSProp, whose update is $\theta_{t+1}=\theta_t-\eta P_t^{-1}\hat m_t$ with a diagonal preconditioner $P_t$ ($\hat m_t$ is the gradient itself for RMSProp). The \term{preconditioned sharpness}\index{preconditioned sharpness} $S_P=\lmax(P_t^{-1}\nabla^2\Loss(\theta_t))$ rises until it reaches the stability threshold of the same method with $P_t$ frozen, which is $38/\eta$ for Adam with $\beta_1=0.9$, and then equilibrates there. The sharpness itself is not held at a threshold: the iterates keep moving into regions of higher curvature, and the preconditioner grows to compensate. Similar effects occur with minibatches, the more so the larger the batch.
\end{observation}

The source is an unrefereed preprint (arXiv 2022; a version appeared at a NeurIPS 2023 workshop), and most of its experiments are full batch, but the claim is easy to check: \cref{fig:eos:network} reproduces it, together with \cref{obs:eos:gd}, in a network small enough to train in half a minute. With minibatches the picture changes.

\begin{observation}[Minibatches; \textcite{JastrzebskiEtAl2020}, \textcite{AndreyevBeneventano2024}]\label[observation]{obs:eos:sgd}
With minibatch stochastic gradient descent the full-batch sharpness levels off below $2/\eta$, at a value set by the step size and the batch size. What settles at $2/\eta$ is a sharpness adapted to the minibatch: the expected curvature of a minibatch loss along its own minibatch gradient, which Andreyev and Beneventano call the \term{batch sharpness}\index{batch sharpness}.
\end{observation}

The first sentence is older and well replicated \parencite{JastrzebskiEtAl2019,JastrzebskiEtAl2020,LeeJang2023}; the second rests on an unrefereed preprint (arXiv 2024, revised in 2025). This chapter is about full batch, and \cref{ch:noise} takes up the noise.

The three observations have one shape: a stability threshold computed from a quadratic model with everything else frozen, and a training process that drives the relevant curvature up to the threshold and then holds it there. Two questions follow. Why does the iteration not diverge once the threshold is crossed, as the quadratic model says it must? And what does training do while it sits at the threshold: where does it go, and what does the oscillation cost? Volume 2 of this series, \emph{The Hydrogen Atoms of Deep Learning}, Chapter 15, answers the first question for a single excursion. In a two-parameter product model, a step size above the threshold makes the loss spike and throws the iterates to a flatter minimizer, a \term{catapult}\index{catapult} in the language of \textcite{LewkowyczEtAl2020}. That model has no progressive sharpening, so after one catapult it stays below the threshold. This chapter adds the sharpening.

\section{The idealization: one oscillating direction}\label{sec:eos:ideal}

The quadratic model freezes the Hessian. To see the curvature respond to the step size we need three ingredients and no more: a direction along which the iterates oscillate when the step is too large, a curvature along that direction that depends on the other parameters, and a training force that pushes the other parameters toward higher curvature. The following family of losses has exactly these.

\begin{assumption}[One oscillating direction]\label[assumption]{ass:eos:onedir}
The parameters split as $\theta=(x,w)\in\R\times\R^{d-1}$, and
\begin{equation}\label{eq:eos:onedir}
  \Loss(x,w)=\Loss_0(w)+\tfrac12\,s(w)\,x^2 ,
\end{equation}
where $\Loss_0$ and $s$ are twice continuously differentiable and $s(w)>\lmax(\nabla^2\Loss_0(w))$ for every $w$.
\end{assumption}

On the floor $x=0$ of the valley the Hessian is block diagonal, with the block $s(w)$ in the direction of $x$ and $\nabla^2\Loss_0(w)$ in the others, so the sharpness there is $s(w)$ and its eigenvector is the $x$-axis. Gradient descent with step size $\eta$ is, exactly,
\begin{equation}\label{eq:eos:gdonedir}
  x_{t+1}=\bigl(1-\eta s(w_t)\bigr)x_t,\qquad
  w_{t+1}=w_t-\eta\nabla\Loss_0(w_t)-\tfrac\eta2\,x_t^2\,\nabla s(w_t).
\end{equation}
The first equation is the quadratic picture: the oscillating coordinate is multiplied by $1-\eta s$, it decays while $\eta s<2$, and above the threshold it flips sign and grows. The second equation contains the new force. The oscillation pushes $w$ down the gradient of the curvature, with a strength proportional to $x^2$: the larger the oscillation, the faster the sharpness falls. Two numbers measure the competition between the training force and the push,
\begin{equation}\label{eq:eos:alphagamma}
  \alpha(w)=-\ip{\nabla\Loss_0(w),\nabla s(w)},\qquad \gamma(w)=\norm{\nabla s(w)}_2^2 .
\end{equation}
Along gradient flow $\dot w=-\nabla\Loss_0(w)$ the sharpness changes at the rate $\dot s=\alpha$, so $\alpha>0$ is progressive sharpening; along the push, a unit of $x^2$ lowers $s$ at the rate $\frac12\gamma$.

The smallest instance of the family linearizes $\Loss_0$ and $s$ near a point at the threshold. If $\Loss_0(w)=\Loss_0(0)+\ip{g,w}$ and $s(w)=2/\eta+\ip{c,w}$ with $c\ne0$, then $y=\ip{c,w}$ is the excess of the curvature over $2/\eta$, and \eqref{eq:eos:gdonedir} gives $x_{t+1}=-(1+\eta y_t)x_t$ and $y_{t+1}=y_t-\eta\ip{c,g}-\frac\eta2\norm c_2^2x_t^2$. The coupling $\norm{c}_2^2$ can be absorbed into $x$: the pair $(\norm c_2x_t,y_t)$ follows the recursion below with $\alpha=-\ip{g,c}$, while the component of $w$ orthogonal to $c$ moves by $-\eta$ times the component of $g$ orthogonal to $c$ at every step and never feels the oscillation. What remains is a two-dimensional map.

\begin{definition}[The cubic model]\label[definition]{def:eos:cubic}
Fix a step size $\eta>0$ and a rate $\alpha>0$. The \term{cubic model}\index{cubic model} is the loss
\[
  \Loss(x,y)=\tfrac12\Bigl(\frac2\eta+y\Bigr)x^2-\alpha y,\qquad (x,y)\in\R^2,
\]
trained by gradient descent with step size $\eta$:
\begin{equation}\label{eq:eos:cubicmap}
  x_{t+1}=-(1+\eta y_t)\,x_t,\qquad y_{t+1}=y_t+\eta\Bigl(\alpha-\tfrac12x_t^2\Bigr).
\end{equation}
\end{definition}

The name refers to the coupling $\frac12yx^2$, which is cubic. The coordinate $x$ is the displacement along the oscillating direction and $y$ the excess of its curvature over $2/\eta$; on the line $x=0$ the sharpness is exactly $2/\eta+y$, and elsewhere it is the largest eigenvalue of $\bigl(\begin{smallmatrix}2/\eta+y&x\\x&0\end{smallmatrix}\bigr)$, which is $2/\eta+y+O(\eta x^2)$. The loss is unbounded below in $y$: it is a model of a region of the landscape in which training raises the curvature at the constant rate $\alpha$, not of a whole landscape. With $x=0$ the curvature grows by $\eta\alpha$ per step forever; with $x\ne0$ it grows only while $x^2<2\alpha$.

\begin{dictionary}
the coordinate $x$, flipping sign at every step & the component of $\theta_t$ along the top eigenvector of the Hessian (of $P^{-1}\nabla^2\Loss$ for Adam), oscillating at period two\\
the other coordinates $w$, and $\Loss_0(w)$ & all other parameters, and the loss at the centre of the oscillation\\
the curvature $s(w)=2/\eta+y$ & the sharpness $S(\theta_t)$, and its distance $y$ from the threshold\\
the rate $\alpha=-\ip{\nabla\Loss_0,\nabla s}$ & the speed of progressive sharpening, $dS/d\tau$ along gradient flow\\
the coupling $\gamma=\norm{\nabla s}_2^2$ & how fast an oscillation lowers the sharpness, through the third derivative of the loss\\
the time $\tau=\eta t$ & the number of steps times the learning rate, the clock in which gradient flow has unit speed\\
the time average of $x^2$ & the variance $\sigma^2$ of the oscillation in a central flow\\
a frozen preconditioner $P$ & Adam's $\sqrt{\hat v_t}+\epsilon$, which moves on the time scale $1/(1-\beta_2)$
\end{dictionary}

\paragraph{What was thrown away.} The oscillating direction is fixed, while in a network the top eigenvector turns; only one direction is at the threshold, while in a network several can be; the loss is exactly quadratic and even in $x$, so there are no terms $x^3$ or $x^4$ along the oscillation; and in the cubic model the rates $\alpha$ and $\gamma$ are constant, so the model cannot leave the edge by itself. Gradients are exact, and there is no momentum and no preconditioner; for quadratics \cref{sec:eos:thresholds} puts both back. The choice keeps exactly the two forces whose balance is the phenomenon, the training force that raises the curvature and the push of the oscillation that lowers it, and it is what a third-order Taylor expansion of a general loss produces near the threshold (\cref{sec:eos:selfstab}).

\paragraph{What one had to notice.} In the quadratic theory an oscillation of growing amplitude is the beginning of divergence. The insight of \textcite{DamianNichaniLee2023} is that it is also the restoring force: through the third derivative of the loss, the oscillation produces the gradient $\frac12x^2\nabla s$, which lowers the curvature that caused it. This is \term{self-stabilization}\index{self-stabilization}. A second, technical observation makes the model solvable: the right variables are the logarithm of the squared amplitude, $A=\log x^2$, which forgets the sign flips, and the excess curvature $y$. In these variables a cycle of oscillation is a closed curve.

\section{The theory: stability thresholds}\label{sec:eos:thresholds}

All four thresholds of the in-the-wild section come from one stability test. A second-order linear recurrence $e_{t+1}=a e_t-be_{t-1}$ with real $a,b$ has the companion matrix $\bigl(\begin{smallmatrix}a&-b\\1&0\end{smallmatrix}\bigr)$, whose eigenvalues are the roots of $z^2-az+b$; its solutions tend to zero from every pair of initial values if and only if both roots lie in the open unit disk, and by the Schur--Cohn (Jury) test this happens if and only if $\abs b<1$ and $\abs a<1+b$ (\cref{lem:toolkit:recurrence}).

\begin{proposition}[Stability thresholds]\label[proposition]{prop:eos:thresholds}
Let $\Loss(\theta)=\frac12(\theta-\theta^\star)^\top H(\theta-\theta^\star)$ with $H$ symmetric positive semidefinite and $S=\lmax(H)>0$, and let $\eta>0$ and $0\le\beta<1$. Each of the following iterations converges to a minimizer of $\Loss$ from every initialization if and only if the stated inequality holds.
\begin{enumerate}
\item Gradient descent, $\theta_{t+1}=\theta_t-\eta\nabla\Loss(\theta_t)$: \ $\eta S<2$.
\item Heavy ball, $\theta_{t+1}=\theta_t-\eta\nabla\Loss(\theta_t)+\beta(\theta_t-\theta_{t-1})$: \ $\eta S<2+2\beta$.
\item Momentum as a moving average, $m_t=\beta m_{t-1}+(1-\beta)\nabla\Loss(\theta_t)$ and $\theta_{t+1}=\theta_t-\eta m_t$: \ $\eta S<(2+2\beta)/(1-\beta)$.
\item Adam with a frozen preconditioner: $m_t$ as in (iii) with $\beta=\beta_1$, and $\theta_{t+1}=\theta_t-\eta P^{-1}m_t$ with a fixed symmetric $P\succ0$: \ $\eta S_P<(2+2\beta_1)/(1-\beta_1)$, where $S_P=\lmax(P^{-1}H)$.
\end{enumerate}
For $\beta_1=0.9$ the threshold in (iv) is $\eta S_P<38$.
\end{proposition}

\begin{proof}
The plan is to split each iteration into scalar recursions, one for each eigenvector of the Hessian, and to apply the stability test to each. Let $Hq_i=\lambda_iq_i$ with orthonormal $q_1,\dots,q_d$, and write $e_{t,i}=\ip{q_i,\theta_t-\theta^\star}$. Since $\nabla\Loss(\theta)=H(\theta-\theta^\star)$, every iteration acts on each $e_{t,i}$ separately, and it converges to a minimizer exactly when every $e_{t,i}$ converges: the minimizers are the points $\theta^\star+v$ with $Hv=0$.

(i) Here $e_{t+1,i}=(1-\eta\lambda_i)e_{t,i}$. If $\lambda_i=0$ the component is constant. If $\lambda_i>0$, the only possible limit is $0$, and $(1-\eta\lambda_i)^te_{0,i}$ tends to $0$ for every $e_{0,i}$ if and only if $\abs{1-\eta\lambda_i}<1$, that is, $\eta\lambda_i<2$. All components converge if and only if this holds for the largest eigenvalue.

(ii) Here $e_{t+1,i}=(1+\beta-\eta\lambda_i)e_{t,i}-\beta e_{t-1,i}$. If $\lambda_i=0$, the roots of $z^2-(1+\beta)z+\beta$ are $1$ and $\beta$, so $e_{t,i}=a+b\beta^t$ converges. If $\lambda_i>0$, a limit $e$ would satisfy $\eta\lambda_ie=0$, so the component converges for every initialization if and only if it tends to zero, and by the stability test this happens if and only if $\beta<1$, which is assumed, and $\abs{1+\beta-\eta\lambda_i}<1+\beta$, which means $0<\eta\lambda_i<2+2\beta$.

(iii) For $t\ge1$, $\theta_{t+1}-\theta_t=-\eta m_t=-\eta\beta m_{t-1}-\eta(1-\beta)\nabla\Loss(\theta_t)=\beta(\theta_t-\theta_{t-1})-\eta(1-\beta)\nabla\Loss(\theta_t)$. This is the heavy ball (ii) with step size $\eta(1-\beta)$, started from $\theta_0$ and $\theta_1=\theta_0-\eta\beta m_{-1}-\eta(1-\beta)\nabla\Loss(\theta_0)$. If $\beta>0$, every $\theta_1$ arises from some $m_{-1}$, so (ii) gives convergence for every initialization if and only if $\eta(1-\beta)S<2+2\beta$; if $\beta=0$, the iteration is (i).

(iv) Put $\varphi_t=P^{1/2}(\theta_t-\theta^\star)$ and $\mu_t=P^{-1/2}m_t$. Multiplying the two updates by $P^{1/2}$ and $P^{-1/2}$ gives $\mu_t=\beta_1\mu_{t-1}+(1-\beta_1)H_P\varphi_t$ and $\varphi_{t+1}=\varphi_t-\eta\mu_t$, with $H_P=P^{-1/2}HP^{-1/2}$. This is iteration (iii) on the quadratic with Hessian $H_P$, which is symmetric positive semidefinite and similar to $P^{-1}H=P^{-1/2}H_PP^{1/2}$, so $\lmax(H_P)=S_P$, and $\varphi\mapsto P^{-1/2}\varphi+\theta^\star$ maps the minimizers of one problem onto those of the other. Finally $(2+1.8)/0.1=38$.
\end{proof}

Three remarks place the proposition. First, the moving-average form (iii), which Adam uses, is the heavy ball with the step size scaled by $1-\beta$, so at a fixed $\eta$ it tolerates a curvature $(1+\beta)/(1-\beta)$ times larger than gradient descent: $19$ times for $\beta=0.9$ and $199$ times for $\beta=0.99$. Second, Adam's bias corrections multiply its step by factors that tend to one, and its preconditioner is not frozen, so (iv) is a statement about a snapshot; \cref{obs:eos:adam} says that the snapshot's threshold is where Adam's preconditioned sharpness equilibrates. Third, a threshold says when the quadratic model of the loss at the current point is unstable, and nothing about what happens next. That is the content of the rest of the chapter. Nesterov's threshold, the third case of \cref{obs:eos:gd}, is \cref{ex:eos:nesterov}.

One more fact about snapshots is needed later. Adam's preconditioner can shrink, and its preconditioned sharpness rise, only at a speed that $\beta_2$ fixes.

\begin{proposition}[How fast the preconditioner can sharpen]\label[proposition]{prop:eos:speed}
Let $v_{t+1}=\beta_2v_t+(1-\beta_2)g_{t+1}^2$ coordinatewise, with $0<\beta_2<1$ and $v_t>0$, let $P_t=\diag(\sqrt{v_t})+\epsilon I$ with $\epsilon\ge0$, and let $H$ be symmetric with $\lmax(P_t^{-1}H)>0$. Then $P_{t+1}\succeq\sqrt{\beta_2}\,P_t$ and
\[
  \lmax(P_{t+1}^{-1}H)\le\beta_2^{-1/2}\,\lmax(P_t^{-1}H).
\]
\end{proposition}

\begin{proof}
Coordinatewise, $v_{t+1}\ge\beta_2v_t$, so $\sqrt{v_{t+1}}+\epsilon\ge\sqrt{\beta_2}\sqrt{v_t}+\epsilon\ge\sqrt{\beta_2}(\sqrt{v_t}+\epsilon)$, which is $P_{t+1}\succeq\sqrt{\beta_2}P_t$ for diagonal matrices. For $P\succ0$ the largest eigenvalue of $P^{-1}H$ is the largest value of $u^\top Hu/u^\top Pu$ over $u\ne0$, and when it is positive it is attained at a $u$ with $u^\top Hu>0$; for such $u$ the quotient for $P_{t+1}$ is at most $\beta_2^{-1/2}$ times the quotient for $P_t$.
\end{proof}

At a fixed Hessian the preconditioned sharpness can therefore grow by at most $\frac12\log(1/\beta_2)\approx\frac12(1-\beta_2)$ per step in logarithm: by a factor $e$ in no fewer than about $2000$ steps when $\beta_2=0.999$. The bound is attained when the new squared gradients are negligible next to $v_t$, which is the case once the iterates have nearly converged. The bias correction $\hat v_t=v_t/(1-\beta_2^{t+1})$ adds a factor $(1-\beta_2^{t+1})/(1-\beta_2^{t+2})$, which matters only for $t\lesssim1/(1-\beta_2)$.

\section{The theory: self-stabilization}\label{sec:eos:selfstab}

\subsection{Where the cubic model comes from}

The cubic model is what a general loss looks like near the threshold, to leading order. The derivation is due to \textcite{DamianNichaniLee2023}, and we give it in the form we need.

\begin{heuristic}[Heuristic derivation of the cubic model]
Let $\Loss$ be three times continuously differentiable and let $\theta^\star$ be a point with $S(\theta^\star)=2/\eta$, a simple top eigenvalue with unit eigenvector $u$, and $\ip{\nabla\Loss(\theta^\star),u}=0$. A simple eigenvalue is differentiable, with derivative $dS=u^\top(dH)u$, so $\nabla S(\theta^\star)=\nabla^3\Loss(\theta^\star)[u,u,\cdot\,]$. Write $\theta=\theta^\star+xu+\delta$ with $\delta\perp u$, put $y=\ip{\nabla S(\theta^\star),\delta}$, and assume $\ip{\nabla S(\theta^\star),u}=0$. The second-order Taylor expansion of the gradient,
\[
  \nabla\Loss(\theta)\approx\nabla\Loss(\theta^\star)+\nabla^2\Loss(\theta^\star)(xu+\delta)+\tfrac12\nabla^3\Loss(\theta^\star)[xu+\delta,xu+\delta,\cdot\,],
\]
has the component $\frac2\eta x+xy$ along $u$, after dropping $\frac12x^2\nabla^3\Loss[u,u,u]$ and $\frac12\nabla^3\Loss[u,\delta,\delta]$, and the component $-\alpha+\frac12\gamma x^2$ along $\nabla S$, with $\alpha=-\ip{\nabla\Loss(\theta^\star),\nabla S(\theta^\star)}$ and $\gamma=\norm{\nabla S(\theta^\star)}_2^2$, after dropping $\ip{\nabla S,\nabla^2\Loss\,\delta}$ and the other terms that involve $\delta$. A step of gradient descent therefore maps $(x,y)$ to approximately $\bigl(-(1+\eta y)x,\;y+\eta(\alpha-\frac12\gamma x^2)\bigr)$, which is \eqref{eq:eos:cubicmap} for $(\sqrt\gamma\,x,y)$. The motion of $\delta$ orthogonal to $\nabla S$ is gradient descent on $\Loss$, undisturbed by the oscillation.
\end{heuristic}

The dropped terms are of two kinds. Some describe how the reference point, its eigenvector and the rates $\alpha$ and $\gamma$ drift; Damian, Nichani and Lee follow them by moving the reference point along a constrained trajectory, and \cref{sec:eos:central} averages over them. Others are higher-order terms along the oscillation itself, such as $x^2\nabla^3\Loss[u,u,u]$ and a fourth-order term; \cref{sec:eos:misses} shows that they decide what happens over long times.

\subsection{Scales, and why two steps}

The cubic map flips the sign of $x$ at every step, so it is not close to the identity, and no differential equation can describe $x_t$ itself. The squared amplitude $x_t^2$ does not flip: by \eqref{eq:eos:cubicmap}, $x_{t+1}^2=(1+\eta y_t)^2x_t^2$, and every quantity that is even in $x$ evolves by a map close to the identity. Equivalently, the two-step map sends $x_t$ to $(1+\eta y_{t+1})(1+\eta y_t)x_t$, without a sign flip. A continuous limit can describe the envelope of the oscillation, not the oscillation.

Which quantities are of order one? In the variables $A=\log x^2$ and $y$, the map \eqref{eq:eos:cubicmap} reads
\begin{equation}\label{eq:eos:Amap}
  A_{t+1}=A_t+2\log(1+\eta y_t),\qquad y_{t+1}=y_t+\eta\Bigl(\alpha-\tfrac12e^{A_t}\Bigr).
\end{equation}
If $\alpha$, $y$ and $e^A=x^2$ are of order one, then each step changes $A$ and $y$ by amounts of order $\eta$, and the natural clock is $\tau=\eta t$, the time of gradient flow, in which the sharpness rises at the rate $\alpha$ when $x=0$. In this scaling the curvature along the oscillation is within $O(1)$ of $2/\eta$, so $\eta S-2=\eta y$ is of order $\eta$, and one cycle lasts a time of order $1/\sqrt\alpha$ in $\tau$, that is, of order $1/(\eta\sqrt\alpha)$ steps. The model has a single dimensionless parameter: with $a=A-\log(2\alpha)$, $Y=y/\sqrt\alpha$ and $s=\sqrt\alpha\,\tau$, the map becomes $a\mapsto a+2\log(1+\eps Y)$, $Y\mapsto Y+\eps(1-e^a)$, with
\[
  \eps=\eta\sqrt\alpha .
\]
Its square $\eps^2=\eta^2\alpha$ is the amount by which progressive sharpening raises $\eta S$ in one step. The continuous limit is $\eps\to0$: sharpening that is slow on the scale of a step, which is the regime of \cref{obs:eos:gd}.

\subsection{The theorem}

The limit of \eqref{eq:eos:Amap} is the differential equation
\begin{equation}\label{eq:eos:flow}
  \frac{dA}{d\tau}=2y,\qquad \frac{dy}{d\tau}=\alpha-\tfrac12e^{A}.
\end{equation}
Write $\phi(A)=\frac12e^A-\alpha A$ and
\begin{equation}\label{eq:eos:H}
  \cH(A,y)=y^2+\phi(A)=y^2+\tfrac12e^A-\alpha A .
\end{equation}
The function $\phi$ is strictly convex, with $\phi''=\frac12e^A>0$, and it tends to $+\infty$ as $A\to\pm\infty$; its minimum is at $A^*=\log(2\alpha)$, where $\phi(A^*)=\alpha(1-\log2\alpha)=:\eosHmin$. So $\cH$ has a unique minimum $\eosHmin$, at $(A^*,0)$, which corresponds to the orbit $x=\pm\sqrt{2\alpha}$, $y=0$ of period two of gradient descent. For $h>\eosHmin$ let $A_-<A^*<A_+$ be the two solutions of $\phi(A)=h$.

\begin{theorem}[Self-stabilization in the continuous limit]\label{thm:eos:selfstab}
Let $\alpha>0$, let $(A(\tau),y(\tau))$ solve \eqref{eq:eos:flow} from a point other than $(A^*,0)$, and let $h=\cH(A(0),y(0))$.
\begin{enumerate}
\item The energy $\cH$ is constant along the solution.
\item The solution is periodic, with period
\[
  T(h)=\int_{A_-}^{A_+}\frac{dA}{\sqrt{h-\phi(A)}}<\infty .
\]
\item Over each period the time averages are $\langle y\rangle=0$ and $\langle e^A\rangle=2\alpha$: the curvature along the oscillation averages exactly $2/\eta$, and the squared amplitude $x^2$ averages exactly $2\alpha$.
\item Gradient descent \eqref{eq:eos:cubicmap} started at a point $(x_0,y_0)$ with $\log x_0^2=A(0)$ and $y_0=y(0)$ follows the solution: for every $\tau_{\max}>0$ there are constants $C$ and $\eta_0>0$, which depend on $\alpha$, $A(0)$, $y(0)$ and $\tau_{\max}$ but not on $\eta$, such that for every $\eta\in(0,\eta_0]$ and every integer $t$ with $0\le t\le\tau_{\max}/\eta$,
\[
  \abs{\log x_t^2-A(\eta t)}+\abs{y_t-y(\eta t)}\le C\eta .
\]
\end{enumerate}
\end{theorem}

The theorem says that, to leading order in $\eps$, an oscillation at the edge is a closed cycle: the curvature rises above $2/\eta$, the oscillation grows and pushes it back down below $2/\eta$, the oscillation dies, and the cycle repeats with the same energy. Part (iii) is exact, not approximate. Over a cycle the curvature spends as much time above $2/\eta$ as it needs to spend below for its average to be exactly the threshold, and the oscillation carries exactly the energy that cancels progressive sharpening: $\frac12\langle x^2\rangle=\alpha$, or $\frac12\gamma\langle x^2\rangle=\alpha$ before the coupling is absorbed. Part (iv) is the continuous limit; we prove it after the rest.

\begin{proof}[Proof of (i)--(iii)]
The plan is to show that $\cH$ is a conserved energy, that its level sets are closed curves without equilibria, and that a periodic solution has the stated averages because the averages of derivatives over a period vanish.

(i) Along a solution, $\frac{d}{d\tau}\cH=2y\,\dot y+\phi'(A)\dot A=2y(\alpha-\frac12e^A)+(\frac12e^A-\alpha)\,2y=0$.

(ii) The sublevel sets of $\cH$ are compact, because $\cH\ge y^2+\eosHmin$ and $\phi$ tends to infinity at both ends, so by (i) the solution stays in a compact set and exists for all $\tau$. Since $h>\eosHmin$, the level set $\{\cH=h\}$ consists of the two arcs $y=\pm\sqrt{h-\phi(A)}$ over $A_-\le A\le A_+$, joined at the turning points $(A_\pm,0)$. The vector field of \eqref{eq:eos:flow} vanishes only at $(A^*,0)$, which is not on this curve. On the upper arc $\dot A=2y=2\sqrt{h-\phi(A)}>0$, so $A$ increases, and the time it takes to go from a point of the arc to $A_+$ is $\int dA/(2\sqrt{h-\phi(A)})$, which is at most $T(h)/2$. This integral is finite. By convexity, $\phi$ lies below its chord from $(A^*,\eosHmin)$ to $(A_+,h)$, so for $A^*\le A<A_+$,
\[
  h-\phi(A)\ge(h-\eosHmin)\frac{A_+-A}{A_+-A^*},
\]
and $\int_{A^*}^{A_+}dA/\sqrt{h-\phi(A)}\le2(A_+-A^*)/\sqrt{h-\eosHmin}$; the same argument bounds the integral over $[A_-,A^*]$. At the turning point $(A_+,0)$ we have $\dot y=-\phi'(A_+)<0$, because $A_+>A^*$, so the solution passes to the lower arc, along which $A$ decreases, reaches $(A_-,0)$ in the time $\int_{A_-}^{A_+}dA/(2\sqrt{h-\phi(A)})$, and there passes to the upper arc because $\dot y=-\phi'(A_-)>0$. A solution that starts on the lower arc does the same in the other order. After the time $T(h)$, the sum of the two transit times, the solution is back at its starting point, and by uniqueness of solutions it is periodic.

(iii) Over a period $[\tau_0,\tau_0+T]$ both $A$ and $y$ return to their initial values. Integrating the two equations of \eqref{eq:eos:flow} over the period gives $0=A(\tau_0+T)-A(\tau_0)=2\int y\,d\tau$ and $0=y(\tau_0+T)-y(\tau_0)=\int(\alpha-\frac12e^A)\,d\tau$, which are the two averages. The curvature along the oscillation is $2/\eta+y$, and $e^A=x^2$.
\end{proof}

The argument for (iii) is the one that gives the averages of the predator--prey cycles of Lotka and Volterra: over a closed cycle the average of a derivative is zero.\sidenote{The analogy is only partial. In the predator--prey system both populations are positive and enter through logarithms; here $y$ changes sign, and \eqref{eq:eos:flow} is the motion of a particle in the potential well $\phi$, since $d^2A/d\tau^2=-2\phi'(A)$.} The small cycles have the period of the linearization: near $(A^*,0)$, $d^2A/d\tau^2=2\alpha-e^A\approx-2\alpha(A-A^*)$, so $T(h)\to2\pi/\sqrt{2\alpha}$ as $h\downarrow\eosHmin$ (\cref{ex:eos:period}), about $2\pi/(\eta\sqrt{2\alpha})$ steps. Larger cycles are slower, because they spend a long time with $A$ very negative, an oscillation so small that nothing happens except that the curvature climbs back at the rate $\alpha$.

Part (iv) is an error estimate of the kind used for the Euler method, with one difference: the map \eqref{eq:eos:Amap} is not the Euler method of \eqref{eq:eos:flow}, because of the logarithm, but it agrees with it to second order in $\eta$.

\begin{lemma}[The continuous limit]\label[lemma]{lem:eos:limit}
Let $z(\tau)=(A(\tau),y(\tau))$ solve \eqref{eq:eos:flow} on $[0,\tau_{\max}]$, and let $K$ be the set of points within $\ell_1$-distance $1$ of the curve $\{z(\tau):0\le\tau\le\tau_{\max}\}$. Let $Y=\max_K\abs y$ and $E=\max_Ke^A$, and put $L=\max\{4,E/2\}$ and $C_1=\alpha+\frac12E+2Y^2+\frac12EY$. Let $z_t=(A_t,y_t)$ follow \eqref{eq:eos:Amap} from $z_0=z(0)$. If $\eta\le1/(2Y)$ and $C_1\eta\,(e^{L\tau_{\max}}-1)\le L$, then for every integer $t$ with $0\le t\le\tau_{\max}/\eta$ the iterate $z_t$ lies in $K$, $x_t\ne0$, and
\[
  \norm{z_t-z(\eta t)}_1\le\frac{C_1}{L}\bigl(e^{L\eta t}-1\bigr)\,\eta .
\]
\end{lemma}

\begin{proof}
The plan is to bound the error made in one step by the exact solution, to show that the map \eqref{eq:eos:Amap} is Lipschitz with constant $1+L\eta$ on $K$, and to add up the errors with the discrete Grönwall inequality (\cref{thm:toolkit:gronwall}). Write $F$ for the map \eqref{eq:eos:Amap}. The set $K$ is compact, and on it $\abs{\eta y}\le\frac12$, so $F$ is defined and smooth there.

\emph{One step of the exact solution.} Along the solution $\abs y\le Y$ and $e^A\le E$, so $\abs{\dot y}\le\alpha+\frac12E$ and $\abs{\dot A}\le2Y$. Fix $\tau\le\tau_{\max}-\eta$ and write $y=y(\tau)$, $A=A(\tau)$. The first component of $z(\tau+\eta)-F(z(\tau))$ is
\[
  \int_0^\eta2\bigl(y(\tau+r)-y\bigr)\,dr+2\bigl(\eta y-\log(1+\eta y)\bigr).
\]
The integral is at most $2\int_0^\eta r(\alpha+\frac12E)\,dr=(\alpha+\frac12E)\eta^2$ in absolute value. For the second term we use $0\le u-\log(1+u)\le u^2$ for $u\ge-\frac12$: the function $u^2-u+\log(1+u)$ vanishes at $0$ and has derivative $u(1+2u)/(1+u)$, which has the sign of $u$ there. So the second term lies in $[0,2\eta^2Y^2]$. The second component is $-\frac12\int_0^\eta(e^{A(\tau+r)}-e^{A})\,dr$, and $\abs{e^{A(\tau+r)}-e^A}\le E\abs{A(\tau+r)-A}\le2EYr$, so it is at most $\frac12EY\eta^2$. Together, $\norm{z(\tau+\eta)-F(z(\tau))}_1\le C_1\eta^2$.

\emph{Lipschitz bound.} Let $z,z'\in K$ with the segment between them in $K$. The derivative of $2\log(1+\eta y)$ in $y$ is $2\eta/(1+\eta y)\le4\eta$, and the derivative of $\frac12\eta e^A$ in $A$ is at most $\frac12\eta E$, so
\[
  \norm{F(z)-F(z')}_1\le\Bigl(1+\tfrac12\eta E\Bigr)\abs{A-A'}+(1+4\eta)\abs{y-y'}\le(1+L\eta)\norm{z-z'}_1 .
\]

\emph{Adding up.} Let $e_t=\norm{z_t-z(\eta t)}_1$ and suppose that $e_s\le1$ for $s\le t$, as holds for $t=0$. Then $z_t$ lies in the $\ell_1$-ball of radius $1$ around $z(\eta t)$, a convex subset of $K$, and the two bounds give $e_{t+1}\le\norm{F(z_t)-F(z(\eta t))}_1+\norm{F(z(\eta t))-z(\eta t+\eta)}_1\le(1+L\eta)e_t+C_1\eta^2$. With $e_0=0$ this gives, by induction,
\[
  e_{t+1}\le C_1\eta^2\sum_{s=0}^{t}(1+L\eta)^s=\frac{C_1\eta}{L}\bigl((1+L\eta)^{t+1}-1\bigr)\le\frac{C_1\eta}L\bigl(e^{L\eta(t+1)}-1\bigr),
\]
which is at most $1$ as long as $\eta(t+1)\le\tau_{\max}$, by the assumption on $\eta$. So the induction continues up to $t=\tau_{\max}/\eta$. Finally, $z_t\in K$ gives $1+\eta y_t\ge\frac12$, so no step sends $x$ to zero.
\end{proof}

\begin{proof}[Proof of \cref{thm:eos:selfstab}(iv)]
Apply \cref{lem:eos:limit} with $C=(C_1/L)(e^{L\tau_{\max}}-1)$ and $\eta_0=\min\{1/(2Y),\,1/C\}$, and note that $\log x_t^2=A_t$ by \eqref{eq:eos:Amap}.
\end{proof}

\subsection{What the discrete map does}

The theorem is about a limit. For the map itself, three identities hold exactly, at every step size.

\begin{lemma}[Exact identities]\label[lemma]{lem:eos:identities}
Let $(x_t,y_t)$ follow \eqref{eq:eos:cubicmap} from $x_0\ne0$, and suppose that $1+\eta y_t>0$ for $0\le t<T$. Write $A_t=\log x_t^2$. Then:
\begin{enumerate}
\item $\displaystyle\frac1T\sum_{t<T}x_t^2=2\alpha-\frac{2(y_T-y_0)}{\eta T}$;
\item $\displaystyle\frac1T\sum_{t<T}\log(1+\eta y_t)=\frac{A_T-A_0}{2T}$, and hence $\displaystyle\frac1T\sum_{t<T}y_t\ge\frac1\eta\bigl(e^{(A_T-A_0)/(2T)}-1\bigr)$;
\item for $t<T$, with $f_t=\alpha-\frac12x_t^2$,
\begin{equation}\label{eq:eos:dH}
  \cH(A_{t+1},y_{t+1})-\cH(A_t,y_t)=\eta^2f_t^2+\tfrac12\eta^2x_t^2y_t^2+2\alpha\bigl(\eta y_t-\log(1+\eta y_t)\bigr)\ \ge\ 0,
\end{equation}
with equality if and only if $x_t^2=2\alpha$ and $y_t=0$.
\end{enumerate}
\end{lemma}

\begin{proof}
By \eqref{eq:eos:cubicmap}, $x_{t+1}^2=(1+\eta y_t)^2x_t^2>0$ for $t<T$, so every $A_t$ with $t\le T$ is defined and \eqref{eq:eos:Amap} holds. Summing the $y$-update over $t<T$ gives $y_T-y_0=\eta\sum_{t<T}(\alpha-\frac12x_t^2)$, which is (i). Summing the $A$-update gives the identity in (ii), and since $\log$ is concave, $\log\bigl(\frac1T\sum_t(1+\eta y_t)\bigr)\ge\frac1T\sum_t\log(1+\eta y_t)$, which is the inequality. For (iii), drop the index $t$, and write $u=\eta y$, so that $e^{A_{t+1}}=(1+u)^2e^A$ and $A_{t+1}-A=2\log(1+u)$. Expanding,
\[
  \cH(A_{t+1},y_{t+1})-\cH(A,y)=2\eta yf+\eta^2f^2+\tfrac12e^A\bigl(2u+u^2\bigr)-2\alpha\log(1+u).
\]
The terms of first order combine: $2\eta yf+e^Au=2u(\alpha-\frac12e^A)+e^Au=2\alpha u$. What remains is $\eta^2f^2+\frac12u^2e^A+2\alpha(u-\log(1+u))$, which is \eqref{eq:eos:dH}. Each term is nonnegative, the last because $u-\log(1+u)\ge0$ for $u>-1$ with equality only at $u=0$. The sum vanishes if and only if $f=0$ and $u=0$, that is, $x^2=2\alpha$ and $y=0$.
\end{proof}

Identity (i) is the discrete form of $\langle x^2\rangle=2\alpha$, and it holds along every orbit with no limit taken: the squared oscillation averages $2\alpha$ up to a boundary term that is small whenever $y$ stays bounded. Identity (ii) says that the geometric mean of $\abs{x_{t+1}/x_t}=\eta S_t-1$, where $S_t=2/\eta+y_t$ is the curvature along the oscillation, is exactly $\abs{x_T/x_0}^{1/T}$; if the oscillation neither grows nor decays, the arithmetic mean of $\eta S_t$ is at least $2$, by the concavity of the logarithm. This is \cref{obs:eos:gd}(ii)'s ``at or just above $2/\eta$''. Identity (iii) is new information: the energy that the flow conserves increases at every step of the map. The cycles of gradient descent are not closed.

\begin{proposition}[Gradient descent spirals outward]\label[proposition]{prop:eos:drift}
Let $\eta>0$ and $\alpha>0$, and let $F$ be the map \eqref{eq:eos:Amap} on the half-plane $\{1+\eta y>0\}$.
\begin{enumerate}
\item The point $p=(A^*,0)$, the orbit $x=\pm\sqrt{2\alpha}$, $y=0$ of period two of \eqref{eq:eos:cubicmap}, is a fixed point of $F$, and the Jacobian of $F$ at $p$ has the eigenvalues $1\pm i\eta\sqrt{2\alpha}$, of modulus $\sqrt{1+2\alpha\eta^2}>1$.
\item Along every orbit of $F$, $\cH$ is nondecreasing, and it increases strictly at every step taken from a point other than $p$.
\item For every compact set $K$ in the half-plane, every orbit that starts in $K$ at a point other than $p$ leaves $K$.
\end{enumerate}
\end{proposition}

\begin{proof}
(i) At $p$ we have $1+\eta y=1$ and $e^A=2\alpha$, so $F(p)=p$. The Jacobian of $F$ is $\bigl(\begin{smallmatrix}1&2\eta/(1+\eta y)\\-\frac12\eta e^A&1\end{smallmatrix}\bigr)$, which at $p$ is $\bigl(\begin{smallmatrix}1&2\eta\\-\eta\alpha&1\end{smallmatrix}\bigr)$, with characteristic polynomial $(1-\lambda)^2+2\alpha\eta^2$. (ii) is \cref{lem:eos:identities}(iii). (iii) Suppose an orbit $z_t$ starts at $z_0\in K\setminus\{p\}$ and stays in $K$. The set $K'=K\cap\{\cH\ge\cH(z_0)\}$ is compact, contains every $z_t$ by (ii), and does not contain $p$, because $\cH(p)=\eosHmin<\cH(z_0)$. The increment $\cH(F(z))-\cH(z)$ is continuous and positive on $K'$, so it is at least some $c>0$ there, and $\cH(z_t)\ge\cH(z_0)+ct$ for all $t$, which contradicts the boundedness of $\cH$ on $K$.
\end{proof}

So the period-two orbit at the edge, the oscillation that exactly balances progressive sharpening, is unstable for every step size, and every other orbit drifts outward to larger cycles. The drift is slow, and its rate can be computed.\sidenote{Explicit time stepping does the same to every oscillator. For $\ddot q=-\omega^2q$ written as a first-order system, one Euler step multiplies the energy by exactly $1+\eta^2\omega^2$. In \eqref{eq:eos:Amap} the two variables are updated from the same old values, as in the Euler method.}

\begin{heuristic}[Heuristic derivation of the area law]
Let $\eosArea(h)$ be the area enclosed by the level curve $\{\cH=h\}$ in the $(A,y)$-plane. Over one cycle at level $h$, gradient descent stays within $O(\eta)$ of the cycle (\cref{lem:eos:limit}) and takes about $T(h)/\eta$ steps, each of which raises $\cH$ by $\eta^2P+O(\eta^3)$, with $P=f^2+y^2(\alpha+\frac12e^A)$, by expanding \eqref{eq:eos:dH} in $\eta$. Replacing the sum by an integral, the level rises by $\eta\oint P\,d\tau$ per cycle. On the cycle $dy=f\,d\tau$ and $dA=2y\,d\tau$, so $\oint P\,d\tau=\oint f\,dy+\frac12\oint(\alpha+\frac12e^A)\,y\,dA$. The cycle runs clockwise ($A$ increases where $y>0$), and Green's theorem turns the two line integrals into $\iint\frac12e^A$ and $\frac12\iint(\alpha+\frac12e^A)$ over the region $\{\cH\le h\}$. In that region $\iint(\alpha-\frac12e^A)\,dA\,dy=0$, because $\alpha-\frac12e^A=-\partial_A\cH$ and $\cH=h$ at both ends of every horizontal chord, so the average of $e^A$ over the region is $2\alpha$, as over the cycle, and $\oint P\,d\tau=2\alpha\eosArea(h)$. Finally $\eosArea(h)=\int_{A_-}^{A_+}2\sqrt{h-\phi(A)}\,dA$ has the derivative $\eosArea'(h)=T(h)$. So per cycle $\eosArea$ grows by $T(h)\cdot2\alpha\eta\eosArea(h)$, which is $2\alpha\eta^2\eosArea$ per step:
\begin{equation}\label{eq:eos:area}
  \eosArea\bigl(\cH(z_t)\bigr)\approx\eosArea\bigl(\cH(z_0)\bigr)\,e^{2\alpha\eta^2t}.\qedhere
\end{equation}
\end{heuristic}

Near the period-two orbit the area law is exact to first order, by \cref{prop:eos:drift}(i): the linearization multiplies the area of a small ellipse by its determinant $1+2\alpha\eta^2$ at every step. In the units of the theorem the drift takes the time $1/(2\alpha\eta)$ in $\tau$ to grow an area by the factor $e$, while a small cycle lasts $2\pi/\sqrt{2\alpha}$; the drift therefore needs about $1/(2\sqrt2\,\pi\eps)$ cycles, roughly $11$ cycles at $\eps=0.01$. It is invisible over a few cycles and dominant over many.\sidenote{The area $\eosArea$ is $2\pi$ times the action of the cycle in the sense of classical mechanics. Actions are conserved under slow changes of the parameters (adiabatic invariance), which is why the slow variation of $\alpha$ along a training run does not by itself change the size of the cycles, while the discretization does.}

\subsection{What Damian, Nichani and Lee prove}

Their theorem concerns gradient descent on a general loss, not the cubic model. Its assumptions include progressive sharpening along the constrained trajectory (their Assumption 1), a gap below the top eigenvalue, with the second eigenvalue of the Hessian below $c/\eta$ for an absolute constant $c<2$ (their Assumption 2), and bounds on the third derivative of the loss and on its Lipschitz constant. The constrained trajectory is projected gradient descent onto the stable set $\{\theta:S(\theta)\le2/\eta,\ \ip{\nabla\Loss(\theta),u(\theta)}=0\}$, with $u(\theta)$ the top eigenvector, and they show that gradient descent can be coupled with it for $\cT=O(\eps^{-1})$ steps, $\eps$ being the small parameter of their analysis, with the deviation captured by a two-dimensional system of the form \eqref{eq:eos:cubicmap} with time-dependent coefficients. The coupling assumes that the oscillation does not become too small: $\min_{t\le\cT}\abs{x_t}\ge c_1\delta$ for a constant $c_1$ and a scale $\delta$ of their analysis. They note that these conditions are not strictly necessary and that a local coupling of this kind is not sufficient to prove global convergence.

Two features of the statement matter here. The horizon is finite, and the lower bound on $\abs{x_t}$ excludes the long stretches of a large cycle in which the oscillation has almost died out. \Cref{prop:eos:drift} and \eqref{eq:eos:area} concern what lies beyond: a time of order $1/(\alpha\eta^2)$ steps, over which the reduced dynamics itself, exact for the cubic model, is not stable. A coupling theorem over a bounded number of cycles cannot decide whether real dynamics at the edge drifts outward, settles onto a stable period-two orbit, or does something else; that is decided by the terms that the reduction drops, and \cref{sec:eos:misses} shows how.

\subsection{Central flows}\label{sec:eos:central}

Average over the cycles. In the family \eqref{eq:eos:onedir}, suppose the curvature $s(w_t)$ has reached $2/\eta$ and the functions $\Loss_0$ and $s$ change little over one cycle. Then, over a cycle, the pair $(\sqrt{\gamma(w)}\,x,\ s(w)-2/\eta)$ follows the cubic model with the local rate $\alpha(w)$, and the theorem gives the cycle averages.

\begin{heuristic}[Heuristic derivation of the central flow]
If $\alpha(w)>0$, the squared oscillation averages $\langle x^2\rangle=2\alpha(w)/\gamma(w)$ over a cycle, by \cref{thm:eos:selfstab}(iii) for the rescaled pair. Averaging the $w$-update of \eqref{eq:eos:gdonedir} over the cycle, and writing $\tau=\eta t$, gives
\begin{equation}\label{eq:eos:central}
  \frac{dw}{d\tau}=-\nabla\Loss_0(w)-\tfrac12\sigma^2(w)\,\nabla s(w),\qquad
  \sigma^2(w)=\frac{2\max\{0,\,-\ip{\nabla\Loss_0(w),\nabla s(w)}\}}{\norm{\nabla s(w)}_2^2},
\end{equation}
while $s(w)=2/\eta$; when $\alpha(w)\le0$ the curvature falls below $2/\eta$, the oscillation decays, and $w$ follows gradient flow. With $\sigma^2=2\alpha/\gamma$, the velocity in \eqref{eq:eos:central} is $-\nabla\Loss_0$ minus its component along $\nabla s$: it is tangent to the surface $\{s=2/\eta\}$. So \eqref{eq:eos:central} is gradient flow projected onto the region $\{s\le2/\eta\}$, and the average of the loss over a cycle is $\Loss_0(w)+\frac12s\langle x^2\rangle=\Loss_0(w)+\sigma^2/\eta$.
\end{heuristic}

\Textcite{CohenEtAl2025} derived the same equation for a general loss and called it a \term{central flow}\index{central flow}: a differential equation for the time-averaged trajectory of gradient descent at the edge of stability, with $\Loss$ and $S$ in place of $\Loss_0$ and $s$. In their words, $\sigma^2$ is the unique value for which the downward force of oscillation-induced sharpness reduction cancels the upward force of progressive sharpening, and the flow projects the direction $\nabla S$ out of the gradient to keep the sharpness at $2/\eta$; they note that this re-derives a main result of \textcite{DamianNichaniLee2023} in a simpler but less rigorous way, and they extend the construction to several eigenvalues at the edge. The status is the same in their paper and here: derived, not proved. Its predictions are nevertheless quantitative: they report that central flows predict long-term trajectories of neural networks with high numerical accuracy. In the cubic model the central flow is the point $y=0$ with $\sigma^2=2\alpha$, which is \cref{thm:eos:selfstab}(iii).

For adaptive optimizers the central flow also moves the preconditioner, and the oscillation acts through it. \Textcite{CohenEtAl2025} derive central flows for RMSProp and Adam, and find what they call acceleration via regularization: the adaptive optimizers implicitly move toward regions of low curvature, in which they can take larger steps.

\section{Reading the theory}\label{sec:eos:reading}

\paragraph{Thresholds in numbers.} \Cref{prop:eos:thresholds} turns the defaults of practice into numbers. For the moving-average form of momentum, the largest stable $\eta S$, or $\eta S_P$ for Adam with a frozen preconditioner, is $(2+2\beta)/(1-\beta)$:
\begin{center}
\small
\begin{tabular}{@{}lccccc@{}}
$\beta$ (or $\beta_1$) & $0$ & $0.5$ & $0.9$ & $0.95$ & $0.99$\\
\hline
largest stable $\eta S$ & $2$ & $6$ & $38$ & $78$ & $398$
\end{tabular}
\end{center}
The left panel of \cref{fig:eos:network} shows full-batch gradient descent on a small network: the sharpness rises, overshoots, and then holds at $2/\eta$, within $1.4\%$ on average over the last $1500$ steps for $\eta=0.4$ and within $0.1\%$ for $\eta=0.25$. The right panel shows Adam on the same network: the preconditioned sharpness climbs to $38/\eta$ and stays there, with averages $0.98$ and $0.95$ of the threshold over the last $1500$ steps and excursions to about $1.03$ times it. The climb is close to the limit of \cref{prop:eos:speed} for a fixed Hessian, because the gradients are already small when it happens: for $\eta=0.003$ the preconditioned sharpness grows by the factor $e^{1.80}$ between steps $200$ and $1975$, where the limit allows $e^{1.67}$, and the slow growth of the Hessian itself makes up the difference. Meanwhile the sharpness of Adam's iterates stays between $7.4$ and $8.4$ from step $300$ on, in both runs, far below any threshold of its own: as \cref{obs:eos:adam} says, what Adam holds at its edge is the preconditioned sharpness.

\begin{figure}[t]
\begin{fullwidth}
\centering
\begin{tikzpicture}
\begin{groupplot}[group style={group size=2 by 1,horizontal sep=17mm},wplot,width=82mm,height=50mm,xmin=0,xmax=4000,xtick={0,1000,2000,3000,4000}]
\nextgroupplot[xlabel={step $t$},ylabel={sharpness $S(\theta_t)$},ymin=0,ymax=10,legend style={at={(0.98,0.04)},anchor=south east,font=\scriptsize},legend columns=2]
\addplot[theory,domain=0:4000,samples=2] {5};
\addplot[theory2,domain=0:4000,samples=2] {8};
\addplot[simulation,mark size=1.0pt] table[x=t,y=gdA] {figures/data/08-eos-network.dat};
\addplot[simulation2,mark size=0.9pt] table[x=t,y=gdB] {figures/data/08-eos-network.dat};
\legend{$2/\eta$\,($\eta=0.4$),$2/\eta$\,($\eta=0.25$),$\eta=0.4$,$\eta=0.25$}
\nextgroupplot[xlabel={step $t$},ylabel={preconditioned sharpness $S_P$},ymode=log,ymin=300,ymax=30000,legend style={at={(0.98,0.04)},anchor=south east,font=\scriptsize},legend columns=2]
\addplot[theory,domain=0:4000,samples=2] {3800};
\addplot[theory2,domain=0:4000,samples=2] {12666.67};
\addplot[simulation,mark size=1.0pt,restrict x to domain=25:4000] table[x=t,y=adA] {figures/data/08-eos-network.dat};
\addplot[simulation2,mark size=0.9pt,restrict x to domain=25:4000] table[x=t,y=adB] {figures/data/08-eos-network.dat};
\legend{$38/\eta$\,($\eta=0.01$),$38/\eta$\,($\eta=0.003$),$\eta=0.01$,$\eta=0.003$}
\end{groupplot}
\end{tikzpicture}
\end{fullwidth}
\caption{The thresholds of \cref{prop:eos:thresholds} (lines) and the measured curvature (points) in a tanh network $8\to32\to32\to1$ ($1377$ parameters) trained full batch on the square loss, with $256$ Gaussian inputs labelled by a random teacher (seeds fixed in \texttt{08-eos-network.py}). Left: gradient descent; the sharpness, estimated every $25$ steps by block power iteration with Hessian--vector products. Right: Adam with $\beta_1=0.9$, $\beta_2=0.999$, $\epsilon=10^{-8}$; the preconditioned sharpness $\lmax(P_t^{-1}H_t)$ of the preconditioner $P_t$ used in the update, from step $25$ on. From step $300$ on, the sharpness of Adam's own iterates stays between $7.4$ and $8.4$ in both runs.}
\label{fig:eos:network}
\end{figure}

\paragraph{The cubic model, cycle by cycle.} \Cref{fig:eos:cubic} compares the theorem with the map. On the left, gradient descent with $\alpha=1$ and $\eta=0.002$ follows the closed level curves of $\cH$ for two cycles each: a small one, started at the edge with $y_0=0.5$, and a large one, started below the edge ($y_0=-0.5$) with a small oscillation $x_0=0.01$, as at the onset of the edge of stability. The large cycle is a sequence of catapults: the curvature climbs to $2/\eta+3.0$, the oscillation grows from $0.01$ to about $5$, the curvature is pushed down to $2/\eta-3.0$, and the oscillation dies back to below $0.01$. The overshoot equals the undershoot because $\cH$ is even in $y$; with $\alpha=0$ the same conservation law turns a single excursion to $2/\eta+y_0$ into a landing at $2/\eta-y_0$, the reflection law that volume 2 of this series, Chapter 15, derives for the catapult. A small oscillation at onset means a large cycle: $\cH\approx y_0^2-\alpha\log x_0^2$ is logarithmic in the initial amplitude. The small cycle lasts $T=4.54$ in $\tau$, close to the $2\pi/\sqrt2=4.44$ of the linearization; the large one lasts $7.56$, that is, $3780$ steps. The second loop of the large cycle lies visibly outside the first: the drift.

In the middle, the running averages of $x^2$ and $y$ along the large cycle converge to $2\alpha=2$ and $0$, with oscillations that decay like $1/\tau$ and that gradient descent shares with the flow. After five periods gradient descent's average of $x^2$ is $2.04$, and identity (i) of \cref{lem:eos:identities} accounts for the difference exactly: it is the boundary term $-2(y_T-y_0)/(\eta T)$, which is not zero because the drifting orbit has fallen slightly out of phase with the flow. On the right, the area law \eqref{eq:eos:area}: three runs with different $(\alpha,\eta)$, started at $\cH=1$, collapse onto the line $\log(\eosArea_t/\eosArea_0)=2\alpha\eta^2t$ over the plotted range, in which the area grows by the factor $e^6$; beyond it, $\eta y$ is no longer small and the expansion in $\eta$ fails.

\begin{figure}[t]
\begin{fullwidth}
\centering
\begin{tikzpicture}
\begin{groupplot}[group style={group size=3 by 1,horizontal sep=14mm},wplot,width=58mm,height=50mm]
\nextgroupplot[xlabel={$A=\log x^2$},ylabel={$y=S-2/\eta$},xmin=-11,xmax=4.5,ymin=-3.6,ymax=3.6]
\addplot[theory] table[x=a2,y=y2] {figures/data/08-eos-cubic-levels.dat};
\addplot[theory] table[x=a1,y=y1] {figures/data/08-eos-cubic-levels.dat};
\addplot[simulation,mark size=0.6pt,unbounded coords=discard] table[x=a2,y=y2] {figures/data/08-eos-cubic-phase.dat};
\addplot[simulation2,mark size=0.6pt,unbounded coords=discard] table[x=a1,y=y1] {figures/data/08-eos-cubic-phase.dat};
\nextgroupplot[xlabel={time $\tau=\eta t$},ylabel={running averages},xmin=0,xmax=38,ymin=-1.2,ymax=4.2,legend style={at={(0.98,0.97)},anchor=north east,font=\scriptsize}]
\addplot[guide,domain=0:38,samples=2] {2};
\addplot[guide,domain=0:38,samples=2] {0};
\addplot[theory] table[x=tau,y=flowx2] {figures/data/08-eos-cubic-avg.dat};
\addplot[theory2] table[x=tau,y=flowy] {figures/data/08-eos-cubic-avg.dat};
\addplot[simulation,mark size=0.7pt] table[x=tau,y=gdx2] {figures/data/08-eos-cubic-avg.dat};
\addplot[simulation2,mark size=0.6pt] table[x=tau,y=gdy] {figures/data/08-eos-cubic-avg.dat};
\legend{,,{flow, $x^2$},{flow, $y$},{GD, $x^2$},{GD, $y$}}
\nextgroupplot[xlabel={$2\alpha\eta^2t$},ylabel={$\log(\eosArea_t/\eosArea_0)$},xmin=0,xmax=6.2,ymin=0,ymax=6.4,legend style={at={(0.03,0.97)},anchor=north west,font=\scriptsize}]
\addplot[theory,domain=0:6.2,samples=2] {x};
\addplot[simulation,mark size=1.1pt] table[x=s1,y=r1] {figures/data/08-eos-cubic-area.dat};
\addplot[simulation2,mark size=1.0pt] table[x=s2,y=r2] {figures/data/08-eos-cubic-area.dat};
\addplot[only marks,mark=triangle*,mark size=1.4pt,color=sea,mark options={fill=sea,draw=white,line width=0.2pt}] table[x=s3,y=r3] {figures/data/08-eos-cubic-area.dat};
\legend{area law,{$\alpha=1$, $\eta=0.01$},{$\alpha=1$, $\eta=0.005$},{$\alpha=\frac12$, $\eta=0.02$}}
\end{groupplot}
\end{tikzpicture}
\end{fullwidth}
\caption{The cubic model: theory (curves) against the gradient-descent map \eqref{eq:eos:cubicmap} (points); no randomness (script \texttt{08-eos-cubic.py}). Left: two cycles each of a small and a large orbit, $\alpha=1$, $\eta=0.002$, every $25$th iterate, on the level curves of $\cH$ through their starting points. Middle: running averages of $x^2$ and of $y$ along the large orbit over five periods, for the flow \eqref{eq:eos:flow} and for gradient descent. Right: the area enclosed by the level curve through the current iterate, for three runs started on $y=0$ at $\cH=1$, against the area law \eqref{eq:eos:area}.}
\label{fig:eos:cubic}
\end{figure}

\paragraph{Adam's edge without sharpening.} The preconditioned sharpness can reach its threshold even when the Hessian does not change at all, because the preconditioner shrinks as the gradients do.

\begin{example}[Adam on a parabola]\label[example]{ex:eos:parabola}
Run Adam with $\beta_1=0.9$, $\beta_2=0.999$ and $\epsilon=0$ on $\Loss(x)=\frac12hx^2$. The iterate converges toward zero, the gradients become small next to $\sqrt{v_t}$, and the preconditioned sharpness $h/P_t$ rises at the limit of \cref{prop:eos:speed}. In a run with $h=1$, $\eta=0.01$ and $x_0=1$ it grows from $\eta S_P=0.066$ at step $1000$ to the threshold $38$ at step $13{,}254$, which is the rate $\frac12\log(1/\beta_2)$ per step corrected for the bias correction of $\hat v$. It then oscillates between about $31$ and $40$, around an average of $35.4$, with $\abs{x}$ of order $0.03\,\eta$. Adam with $\epsilon=0$ is invariant under rescaling the loss, and rescaling $x$ rescales $\eta$; so the range and the average of $\eta S_P$ do not depend on $h$ or on $\eta$, which the runs confirm.
\end{example}

The example separates two mechanisms that the network run mixes. Adam can sharpen through its preconditioner alone, at a speed that $\beta_2$ fixes, and its self-stabilization then runs through $v_t$ rather than through the third derivative of the loss: a growing oscillation raises the squared gradients, $v_t$ grows, and $S_P$ falls below the threshold. The average sits about $7\%$ below the threshold, unlike the cubic model's exact $2/\eta$, because the oscillation grows fast above the threshold but decays slowly below it: the two roots of the momentum recursion have the product $\beta_1$, so the amplitude is multiplied by at least $\sqrt{\beta_1}$ per step however far below the threshold $\eta S_P$ is, and the cycle spends most of its time below (\cref{ex:eos:adam1d}).

\paragraph{Where the agreement must fail.} The theorem is a statement about $\eps\to0$ with the energy fixed. Large cycles fail first: their curvature excursion $\sqrt{h-\eosHmin}$ grows with the energy, and once $\eta y$ is no longer small the expansion behind \eqref{eq:eos:area} breaks down. In the run that the right panel of \cref{fig:eos:cubic} does not show, with $\eta=0.01$, $\alpha=1$ and $\cH_0=10$, the drift carries the orbit to cycles in which $x$ falls below $10^{-300}$; in floating point $x$ underflows to zero after about $4\cdot10^4$ steps, after which the iterates never oscillate again and the curvature grows without bound. In exact arithmetic the cycles keep growing until one of them bursts to an oscillation of astronomical size. Neither is what networks do, and \cref{sec:eos:misses} says why.

\section{What the theory misses}\label{sec:eos:misses}

\paragraph{What was idealized away.} The rate of progressive sharpening $\alpha$ is an input, constant in the cubic model and a given function of $w$ in \eqref{eq:eos:onedir}; why it is positive in networks is not addressed. One direction is at the edge, with a fixed eigenvector and a spectral gap below it, while in networks the top eigenvector turns and several eigenvalues can reach $2/\eta$ together. The loss is exactly quadratic and even in the oscillating coordinate: no terms $x^3$ or $x^4$, which are precisely the terms that decide the long-time behaviour below. The gradients are exact. And the thresholds for momentum and Adam are proved for quadratics with a frozen preconditioner, while the self-stabilization theory is proved for plain gradient descent only.

\paragraph{Where it breaks.} (1) \emph{Long times.} \Cref{prop:eos:drift} shows that the cubic model has no stable cycle: its period-two orbit repels at the rate $\alpha\eta^2$ per step. A fourth-order term $\frac1{24}\kappa x^4$ added to the loss changes the continuous limit at order one: the equations become $dA/d\tau=2y+\frac13\kappa e^A$ and $dy/d\tau=\alpha-\frac12e^A$, and for $\kappa<0$, a loss that grows more slowly than quadratically along the oscillation, the cycles spiral in to a stable period-two orbit with $y=-\kappa\alpha/3>0$, a sharpness slightly above $2/\eta$ (\cref{ex:eos:quartic}). Near the orbit the area rate becomes $2\alpha(\eta+\kappa/3)$ per unit of $\tau$, so the orbit is stable when $\kappa<-3\eta$, a prediction that the exact map confirms. A cubic term $\frac16\kappa_3x^3$ contributes $-\frac12\eta\kappa_3^2e^A$ to $dA/d\tau$, which also damps. In one dimension the combination is the condition $3\Loss'''^2-\Loss''\Loss''''>0$ for a stable period-two orbit near the threshold (\cref{ex:eos:quartic}); \textcite{ChenBruna2023} analyse such orbits through the two-step map. Whether a network's cycles are damped is therefore a property of its fourth derivative along the top eigenvector, which the cubic model sets to zero; the sustained edge of \cref{obs:eos:gd}, hovering ``at or just above'' $2/\eta$, is what damping predicts. (2) \emph{Progressive sharpening.} The model does not say why the sharpness rises, and with the cross-entropy loss it eventually falls \parencite{CohenEtAl2021}; for logistic regression on separable data, gradient descent converges at every step size after a phase in which the loss oscillates \parencite{WuBravermanLee2023,WuEtAl2024}, so there the edge is left for good. (3) \emph{Several eigenvalues.} When two eigenvalues are at the edge, the averaging identity of \cref{lem:eos:identities}(i) becomes a linear system for the two oscillation energies, which can have a negative solution; then one direction cannot stay at the edge (\cref{ex:eos:two}). The single-eigenvalue picture, and the eigengap assumption of \textcite{DamianNichaniLee2023}, fail exactly there. (4) \emph{Noise.} With minibatches the full-batch sharpness stays below $2/\eta$ (\cref{obs:eos:sgd}); the model has no noise and predicts the full-batch threshold. (5) \emph{Adam's moving preconditioner.} \Cref{ex:eos:parabola} shows that Adam reaches its edge and stabilizes there through $v_t$ alone, with an average $7\%$ below the frozen-preconditioner threshold; the cubic mechanism and this one act together in a network, and the theory of this chapter covers neither their interaction nor the observation that Adam keeps moving into regions of higher sharpness \parencite{CohenEtAl2022}.

\paragraph{Repairs and open problems.} Central flows \parencite{CohenEtAl2025} handle several eigenvalues at the edge and adaptive optimizers, and predict long training trajectories, at the level of a derivation. Rod flows \parencite{RegisChewi2026} model gradient descent as an extended object with a centre and an extent; their authors report that rod flow matches central flows on representative networks, contains central flows as a Taylor approximation, and exhibits self-stabilization in quartic potentials, and a second paper extends the construction to Adam, heavy ball, Nesterov and RMSProp \parencite{RegisChewi2026Adam}; both are unrefereed preprints of 2026. Minimalist models prove edge-of-stability behaviour for all time in low dimension: a two-parameter example \parencite{ZhuEtAl2023}, the learning of threshold neurons \parencite{AhnEtAl2023}, normalized gradient descent and a flow that reduces sharpness \parencite{AroraLiPanigrahi2022}, scalar networks \parencite{KreislerEtAl2023}, and a bifurcation analysis that aligns trajectories \parencite{SongYun2023}; \textcite{KalraHeBarkeshli2025} analyse the fixed points of a low-dimensional model and its route to chaos. Progressive sharpening has partial explanations, through the growth of the output layer \parencite{LiWangLi2022} and in second-order regression models \parencite{AgarwalaPedregosaPennington2023}. For minibatches, the batch sharpness of \textcite{AndreyevBeneventano2024} and the sharpness measure of \textcite{LeeJang2023} replace $S$. Open: a model in which progressive sharpening is derived and the long phase at the edge is proved; a global-in-time theorem for gradient descent at the edge beyond low-dimensional examples; the sign of the damping in networks, which can be measured (\cref{ex:eos:drift}); and a theory of Adam's edge with the preconditioner moving.

\section{Back in the wild}\label{sec:eos:back}

\fieldnotes{Learning rates chosen as large as training tolerates; full-batch gradient descent then trains with the sharpness at $2/\eta$, Adam with the preconditioned sharpness at $38/\eta$.}
{One oscillating direction whose curvature depends on the other parameters; its affine case, the cubic model $\frac12(2/\eta+y)x^2-\alpha y$.}
{Thresholds $2/\eta$, $(2+2\beta)/\eta$, $(2+2\beta)/((1-\beta)\eta)$, and $38/\eta$ for Adam with a frozen preconditioner. In the continuous limit, closed cycles with $\langle y\rangle=0$ and $\langle x^2\rangle=2\alpha$; for the map, exact averaging identities and an energy that increases at every step.}
{The sharpness averages $2/\eta$, or slightly more; the oscillation carries $\langle x^2\rangle=2\alpha/\norm{\nabla S}^2$; small cycles last $2\pi/(\eta\sqrt{2\alpha})$ steps; cycles grow by the factor $e$ in area every $1/(2\alpha\eta^2)$ steps unless higher-order terms damp them; the central flow (derived).}
{Why the sharpness rises; several eigenvalues at the edge; noise; a moving preconditioner; the fourth-order terms that decide whether the cycles are damped.}
{The step size sets the curvature, not the reverse (observed; explained in the model). Compare learning rates across momentum values through $(2+2\beta)/(1-\beta)$ (proved for quadratics). The oscillation costs about $\sigma^2/\eta$ of loss, which averaging two consecutive iterates removes (derived in the model; suggested).}

For a practitioner the first lesson is a change of viewpoint. The learning rate is not chosen to respect a curvature; the curvature adapts to the learning rate, rising until the step is barely stable and then staying there. A sweep that keeps the largest learning rate that does not diverge is therefore not choosing a point safely inside a fixed stability region: it is choosing the curvature at which training will run. Divergence happens when self-stabilization fails, for example when the step size grows faster than the curvature can follow, which is the subject of warmup (\cref{ch:warmup}), or when one cycle grows too large, which is one of the anatomies of a loss spike (\cref{ch:spikes}). The second lesson is quantitative but modest. The thresholds tell how momentum moves the edge: at $\beta=0.9$ the moving-average form tolerates a curvature nineteen times larger than plain gradient descent at the same $\eta$, so a learning rate tuned for one $\beta$ is not the right one for another; \cref{ch:momentum} proves the regions and their noise counterparts. And the oscillation is not a malfunction: it carries exactly the energy that progressive sharpening requires, and its cost in loss, about $\sigma^2/\eta$ in the model, is paid by the iterates and not by their average. In the model, the average of two consecutive iterates, $\frac12(x_t+x_{t+1})=-\frac12\eta y_tx_t$, has almost no oscillation, so its loss is lower by nearly $\sigma^2/\eta$; this is a reason to expect averaging of iterates, discussed in \cref{ch:momentum,ch:schedules}, to help at the edge, but it is a suggestion, not a result about networks.

What the theory does not tell a practitioner is equally clear. It does not say why the sharpness rises, so it cannot predict how long a run takes to reach the edge or where it goes along it, beyond the central flow's projection. With noise the edge moves (\cref{ch:noise}), and with Adam the edge involves the preconditioner, whose own dynamics the frozen threshold ignores; why Adam beats gradient descent on transformers is a separate question with several answers (\cref{ch:transformers}).

\begin{exercises}
\item \consolidation (Nesterov's threshold.) Nesterov's momentum is $\theta_{t+1}=\vartheta_t-\eta\nabla\Loss(\vartheta_t)$ with $\vartheta_t=\theta_t+\beta(\theta_t-\theta_{t-1})$. (a) On a quadratic with curvature $\lambda>0$ along an eigenvector, show that the component obeys $e_{t+1}=(1+\beta)(1-\eta\lambda)e_t-\beta(1-\eta\lambda)e_{t-1}$. (b) With the stability test, show that it is stable for every initialization if and only if $0<\eta\lambda<(2+2\beta)/(1+2\beta)$. (c) Compare with the heavy ball for $\beta=0.9$, and with the maximal stable sharpness of Nesterov's method reported by \textcite{CohenEtAl2021}.\label{ex:eos:nesterov}

\item \consolidation (Small cycles.) (a) Show that the linearization of \eqref{eq:eos:flow} at $(A^*,0)$ has the angular frequency $\sqrt{2\alpha}$. (b) Substitute $A=A^*+\sqrt{2(h-\eosHmin)/\alpha}\,\sin\vartheta$ in the formula for $T(h)$ and show that $T(h)\to2\pi/\sqrt{2\alpha}$ as $h\downarrow\eosHmin$. (c) For $\alpha=1$, compute $T$ numerically at $h=\cH(A^*,0.5)$ and compare with the value $4.54$ quoted in \cref{sec:eos:reading}.\label{ex:eos:period}

\item \consolidation (The reduction.) In the family \eqref{eq:eos:onedir} with $\Loss_0(w)=\Loss_0(0)+\ip{g,w}$ and $s(w)=2/\eta+\ip{c,w}$, show that $(\norm c_2x_t,\ip{c,w_t})$ follows \eqref{eq:eos:cubicmap} with $\alpha=-\ip{g,c}$, and that the component of $w_t$ orthogonal to $c$ performs gradient descent on $\Loss_0$ restricted to that subspace. Deduce that the central flow \eqref{eq:eos:central} is exact for this family, in the sense that $\langle y\rangle=0$ and the orthogonal component moves at the velocity $-\eta$ times the projected $g$ per step.\label{ex:eos:reduction}

\item \consolidation (Hellmann--Feynman.) Let $H(\theta)$ be a smooth family of symmetric matrices whose top eigenvalue $S(\theta)$ is simple, with unit eigenvector $u(\theta)$. Differentiate $Hu=Su$ and $u^\top u=1$ to show that $dS=u^\top(dH)\,u$, so that $\nabla S=\nabla^3\Loss[u,u,\cdot\,]$ when $H=\nabla^2\Loss$. Why does the formula fail when the top eigenvalue is not simple?\label{ex:eos:hellmann}

\item \consolidation (Two consecutive iterates.) In the cubic model, show that $\Loss(\frac12(x_t+x_{t+1}),y_t)=\frac18\eta^2y_t^2x_t^2(2/\eta+y_t)-\alpha y_t$, and compare with $\Loss(x_t,y_t)$. Averaging over a cycle in the continuous limit, show that the loss of the iterates exceeds that of the averages of consecutive iterates by $2\alpha/\eta+O(1)$.\label{ex:eos:average}

\item \variation (A quartic term.) Add $\frac1{24}\kappa x^4$ to the cubic model. (a) Show that gradient descent becomes $x_{t+1}=-(1+\eta y_t)x_t-\frac16\eta\kappa x_t^3$, with the same $y$-update. (b) Derive the continuous limit $dA/d\tau=2y+\frac13\kappa e^A$, $dy/d\tau=\alpha-\frac12e^A$, and show that its fixed point is $(\log2\alpha,-\kappa\alpha/3)$, a stable focus if $\kappa<0$ and an unstable one if $\kappa>0$. (c) Show that $d\cH/d\tau=\frac13\kappa e^A(\frac12e^A-\alpha)$, and use $\langle e^A\rangle=2\alpha$ and Jensen's inequality to show that for $\kappa<0$ the energy decreases on average over a cycle of the unperturbed flow. (d) Combine with the area law to predict stability of the period-two orbit of the map when $\kappa<-3\eta$, and check numerically with $\eta=0.01$, $\alpha=1$ and $\kappa\in\{-0.01,-0.05\}$. (e) For a one-dimensional loss with $\Loss''=2/\eta$ at a minimum, show that the two-step map of gradient descent is $x\mapsto x-2(k_2^2+k_3)x^3+O(x^4)$ with $k_2=-\frac12\eta\Loss'''$ and $k_3=-\frac16\eta\Loss''''$, so that the period-two orbit born at the threshold is stable when $3\Loss'''^2-\Loss''\Loss''''>0$.\label{ex:eos:quartic}

\item \variation (Momentum does not change the energy.) Run the heavy ball, $\theta_{t+1}=\theta_t-\eta\nabla\Loss(\theta_t)+\beta(\theta_t-\theta_{t-1})$, on the cubic model with $2/\eta$ replaced by $(2+2\beta)/\eta$. Show that $\frac1T\sum_{t<T}x_t^2=2\alpha-2\bigl((y_T-y_0)-\beta(y_{T-1}-y_{-1})\bigr)/(\eta T)$, so that $\langle x^2\rangle=2\alpha$ along every orbit on which $y$ stays bounded, and derive the same for the moving-average form. What does change with momentum?\label{ex:eos:momentum}

\item \variation (Two directions at the edge.) Let $\Loss=\frac12(2/\eta+y_1)x_1^2+\frac12(2/\eta+y_2)x_2^2-\alpha_1y_1-\alpha_2y_2$ arise, as in \cref{ex:eos:reduction}, from $s_j(w)=2/\eta+\ip{c_j,w}$ and $\Loss_0(w)=\ip{g,w}$, so that $y_j=\ip{c_j,w}$ and $\alpha_j=-\ip{g,c_j}$. (a) Show that $y_{j,t+1}=y_{j,t}+\eta\bigl(\alpha_j-\frac12\sum_kG_{jk}x_{k,t}^2\bigr)$ with $G_{jk}=\ip{c_j,c_k}$. (b) Show that if both oscillations stay bounded and $y$ stays bounded, then the time averages satisfy $G\langle x^2\rangle=2\alpha$. (c) Take unit vectors with $\ip{c_1,c_2}=0.8$ and $(\alpha_1,\alpha_2)=(1,0.5)$, for which $G\xi=2\alpha$ has the solution $\xi=(10/3,-5/3)$. With $\eta=0.002$, starting below the edge with small oscillations, show numerically that the second oscillation dies out while the first stays at the edge with $\langle x_1^2\rangle\approx2$; and that with $(\alpha_1,\alpha_2)=(1,0.9)$ both persist, with averages near $\xi=(1.56,0.56)$.\label{ex:eos:two}

\item \simulation (The area law.) Iterate \eqref{eq:eos:cubicmap} from $y_0=0$ and $\cH_0=1$ for $(\alpha,\eta)\in\{(1,0.01),(1,0.005),(\frac12,0.02)\}$, compute the enclosed area $\eosArea$ of the level curve through each iterate by quadrature, and reproduce the right panel of \cref{fig:eos:cubic}. Then start at $\cH_0=10$ with $\eta=0.01$ and $\alpha=1$, and observe that $x$ underflows to zero in floating point after about $4\cdot10^4$ steps. Repeat in the variables $(A,y)$ of \eqref{eq:eos:Amap}, which cannot underflow, and describe what happens to the cycles.\label{ex:eos:area}

\item \simulation (Adam on a parabola.) Reproduce \cref{ex:eos:parabola}. (a) Measure the rate at which $\log S_P$ grows before the threshold and compare it with \cref{prop:eos:speed}, including the bias correction. (b) Repeat with $\beta_2=0.99$ and with $\beta_1=0.95$, whose threshold is $78$, and report the average and range of $\eta S_P$. (c) Explain why the average lies below the threshold: compare the largest rate at which the oscillation can grow when $\eta S_P$ exceeds $(2+2\beta_1)/(1-\beta_1)$ by a few percent with the rate at which it decays when $\eta S_P$ is $20\%$ below.\label{ex:eos:adam1d}

\item \simulation (Momentum in the network.) Modify \texttt{08-eos-network.py} to run Adam with $\eta=0.01$ and $\beta_1\in\{0.8,0.95\}$, and gradient descent with the moving-average momentum of \cref{prop:eos:thresholds}(iii) at $\beta=0.5$ and $\eta=1$. Check that the preconditioned sharpness settles near $18/\eta$ and somewhat below $78/\eta$, and that the sharpness of the momentum run settles at $6$. Then run the moving average with $\beta=0.9$ and $\eta=4.5$, whose threshold $38/\eta\approx8.4$ is three times the initial sharpness $2.7$, and observe that training diverges within twenty steps. Which hypothesis of \cref{prop:eos:thresholds} fails in the first steps? (\Cref{ch:warmup} takes up the question.)\label{ex:eos:network}

\item \research\hard (Is the edge damped?) For a network at the edge of stability, estimate the third and fourth derivatives of the loss along the top eigenvector, $\nabla^3\Loss[u,u,u]$ and $\nabla^4\Loss[u,u,u,u]$, by finite differences of Hessian--vector products, and compute the damping coefficient of \cref{ex:eos:quartic}(e) together with the discretization's anti-damping $2\alpha\eta$, with $\alpha$ estimated from the sharpening rate before the edge. Does the sign of the total predict whether the oscillation settles onto a period-two orbit? Whether the cycles of networks are damped, and by which terms, is open.\hint{the effective coefficient must include the curvature of $S$ along $u$, since the fourth-order term enters through $\frac{d}{dx}$ of the curvature along the oscillation.}\label{ex:eos:drift}

\item \research (Sharpening as an output.) The cubic model is given progressive sharpening as the rate $\alpha$. Find a model in which $\alpha$ is derived from the data and the architecture, and in which the long phase at the edge, with a decreasing loss, can be proved. \Textcite{AgarwalaPedregosaPennington2023} derive progressive sharpening in second-order regression models; what is open is a model in which both the rate and the long phase are solved.\label{ex:eos:sharpening}
\end{exercises}

\notesandrefs
The threshold $2/\eta$ for gradient descent and the analysis of the heavy ball on quadratics are classical \parencite{Polyak1987}. That stability selects the minima that gradient descent can reach was proposed by \textcite{WuMaE2018}; \textcite{JastrzebskiEtAl2019,JastrzebskiEtAl2020} observed that early in training the sharpness adapts to the step size and the batch size; and \textcite{LewkowyczEtAl2020} described the catapult, solved in volume 2 of this series. \Textcite{CohenEtAl2021} named progressive sharpening and the edge of stability, documented both across architectures and losses, and derived the momentum thresholds; \textcite{CohenEtAl2022} found the adaptive edge of stability.

\Textcite{DamianNichaniLee2023} derived the cubic model by a third-order expansion, identified self-stabilization as the mechanism, and proved the coupling with the constrained trajectory discussed in \cref{sec:eos:selfstab}; an earlier version of their paper appeared at a NeurIPS 2022 workshop. \Textcite{CohenEtAl2025} derived central flows, for gradient descent and for adaptive optimizers. Our \cref{thm:eos:selfstab} is the continuous-limit analysis of the cubic model with the proofs written out. We know no reference for the exact identities of \cref{lem:eos:identities}, for the monotone increase of the energy in \cref{prop:eos:drift}, or for the area law \eqref{eq:eos:area}; they are elementary consequences of the model. Stable period-two orbits near the threshold are the flip bifurcation of dynamical systems \parencite[\S4.4--4.5]{Kuznetsov2004}, and \textcite{ChenBruna2023} study them through the two-step map of gradient descent.

Models with complete proofs of edge-of-stability behaviour are listed in \cref{sec:eos:misses}. Large step sizes for logistic regression, where the edge is eventually left, are analysed by \textcite{WuBravermanLee2023} and \textcite{WuEtAl2024}. Rod flows \parencite{RegisChewi2026,RegisChewi2026Adam} are a recent alternative to central flows.

\paragraph{Translating the literature.} \Textcite{DamianNichaniLee2023} write $x_t$ and $y_t$ as we do, for the displacement along the top eigenvector and its projection on $\nabla S$, and $\alpha=-\ip{\nabla\Loss,\nabla S}$ as we do; their $\beta=\norm{\nabla S}^2$ is our $\gamma$, renamed because $\beta$ is the momentum in this book; their constrained trajectory $\theta_t^\dagger$ is the discrete counterpart of the central flow. \Textcite{CohenEtAl2025} write $w$ for the parameters and count time in steps, so their flows carry a factor $\eta$ that our $\tau=\eta t$ absorbs; their $\sigma^2(t)$ is our $\sigma^2(w(\tau))$, the average of $x^2$ over a cycle. \Textcite{CohenEtAl2021} write $\mathrm{MSS}(\eta,\beta)$ for the maximal stable sharpnesses of \cref{obs:eos:gd}(iii). In \cref{fig:eos:network} we include Adam's bias corrections in the preconditioner, $P_t=(1-\beta_1^{t+1})(\sqrt{\hat v_t}+\epsilon)$, which matters only in the first few hundred steps. Volume 2 of this series, Chapter 15, writes the reduced dynamics as $dA/dt=2\eps A$, $d\eps/dt=\gamma-\nu A$ in steps; its $(A,\eps,\gamma,\nu)$ are our $(x^2,\eta y,\eta^2\alpha,\frac12\eta^2\gamma)$, so its $A$ is our $e^A$.

\printchapterbib
\end{refsection}
